\documentclass[11pt,a4paper]{article}
\usepackage[T1]{fontenc}
\usepackage{lmodern,amsmath,amssymb,amsthm,mathtools,microtype}
\usepackage[margin=27mm]{geometry}
\usepackage[colorlinks=true,linkcolor=blue,citecolor=blue,urlcolor=blue]{hyperref}
\setlength{\emergencystretch}{2em}
\newtheorem{theorem}{Theorem}[section]
\newtheorem{lemma}[theorem]{Lemma}
\newtheorem{proposition}[theorem]{Proposition}
\newtheorem{corollary}[theorem]{Corollary}
\newcommand{\Q}{\mathbb Q}\newcommand{\Pp}{\mathbb P}
\newcommand{\OO}{\mathcal O}
\title{K3 surfaces in the fixed-input pole problem}
\author{}\date{16 November 2026}
\begin{document}\maketitle
\begin{abstract}
For every fixed nonzero pair of coefficients, the normalized one-pole
surface, with its conic parameter still variable, is birational to a
double cover of $\Pp^1\times\Pp^1$ with four singularities of type
$A_3$. Its minimal resolution is a K3 surface. For a finite list of
fixed detector inputs, the component of the common-parameter fiber
product that dominates the parameter line therefore has no rationally
connected smooth projective model. This is a restriction on direct
applications of rational-connectedness methods to the specified
fixed-input model. It is not an obstruction to rational witnesses,
common rational parameters, or constructions that vary the detector
inputs.
\end{abstract}

\section{The projective double cover}

Let $k$ be a field of characteristic zero and fix $C,D\in k^\times$.
Consider the surface
\begin{equation}\label{eq:affine}
 q^2=D(s^2-eC)(s^2+3eC),\qquad
 e=\frac{\tau^2-D}{\tau^2+3D},
\end{equation}
on the open set where the denominator is nonzero. For the normalized
pole of \cite{Guards}, the coefficients are $D=17d$ and
$C=17(17-b)$, with $b=dt^2$ a prepared input. In particular both
coefficients are nonzero.

Put $z=q(\tau^2+3D)$. On
$Q=\Pp^1_{[T:U]}\times\Pp^1_{[S:V]}$, define
\begin{equation}\label{eq:branch}
\begin{aligned}
 F_1&=(T^2+3DU^2)S^2-C(T^2-DU^2)V^2,\\
 F_2&=(T^2+3DU^2)S^2+3C(T^2-DU^2)V^2.
\end{aligned}
\end{equation}
The product $DF_1F_2$ is a section of $\OO_Q(4,4)$.
Let $Y_{C,D}$ be the double cover of $Q$ defined by this section
and the square-root line bundle $L=\OO_Q(2,2)$. Its local equation
on $U=V=1$ is $z^2=DF_1F_2$, so it is birational to
\eqref{eq:affine}.

\begin{theorem}\label{thm:k3}
Over an algebraic closure of $k$, the surface $Y_{C,D}$ is normal
and has exactly four singularities, all of type $A_3$. Its smooth
minimal resolution $\widetilde Y_{C,D}$ is a K3 surface. More
precisely,
\[
 \omega_{\widetilde Y_{C,D}}\simeq\OO_{\widetilde Y_{C,D}},\qquad
 H^1(\widetilde Y_{C,D},\OO)=0.
\]
\end{theorem}

\begin{proof}
All geometric assertions may be checked after passing to an algebraic
closure. First the two curves $B_i=(F_i=0)$ are smooth of bidegree
$(2,2)$. An elementary uniform test is useful here. A biquadratic
form with affine expression
\[
                 a\tau^2s^2+b\tau^2+cs^2+d
\]
is smooth on $\Pp^1\times\Pp^1$ if $a,b,c,d$ and $ad-bc$ are
all nonzero. In the affine chart, a singularity with $\tau s=0$
would force $d=0$ or one of $b,c$ to vanish. Otherwise the two
partial derivatives give $as^2=-b$ and $a\tau^2=-c$; substitution
gives $ad-bc=0$. The three other charts merely permute the four
coefficients. For $F_1$ and $F_2$ the coefficient lists and the
last determinant are
\[
 (1,-C,3D,CD),\quad 4CD;\qquad
 (1,3C,3D,-3CD),\quad -12CD.
\]
Thus both curves are smooth. They are geometrically integral:
their projections to the $\tau$ line are connected double covers
branched at $\tau^2=D$ and $\tau^2=-3D$, four distinct points.
Each consequently has genus one.

The identities
\begin{equation}\label{eq:intersections}
 F_2-F_1=4C(T^2-DU^2)V^2,\qquad
 3F_1+F_2=4(T^2+3DU^2)S^2
\end{equation}
show that their intersections are precisely
\begin{equation}\label{eq:four}
 s=0,\ \tau^2=D;\qquad
 s=\infty,\ \tau^2=-3D.
\end{equation}
There are no intersections over $U=0$. At each of the four
points, the branches have intersection multiplicity two.
Here are exact local normal forms which also determine the double
cover singularities.

Near a point with $s=0$ and $\tau^2=D$, the coordinate
\[
                    x=C\frac{\tau^2-D}{\tau^2+3D}
\]
has nonzero derivative with respect to $\tau$ and is therefore
a local parameter on the first factor. Dividing the ordinate by
the relevant unit and a constant square root, the covering equation
becomes
\begin{equation}\label{eq:normalzero}
 w^2=(s^2-x)(s^2+3x)
     =\frac43s^4-3\left(x-\frac{s^2}{3}\right)^2.
\end{equation}
Near a point with $s=\infty$ and $\tau^2=-3D$, use
$v=1/s$ and $x=(\tau^2+3D)/(\tau^2-D)$. The analogous equation is
\begin{equation}\label{eq:normalinfinity}
 w^2=(x-Cv^2)(x+3Cv^2)=(x+Cv^2)^2-4C^2v^4.
\end{equation}
Both are equivalent over the algebraic closure to
$ab=c^4$, the $A_3$ singularity. Their branch equations consist
of two smooth branches with tangency order two. Away from these
four points the branch divisor is smooth and the double cover is
smooth. The reduced branch divisor also makes the cover integral;
an integral hypersurface with these isolated singularities is normal.

Let $\pi:Y_{C,D}\to Q$ be the covering map. Its structure algebra
and canonical sheaf are
\[
 \pi_*\OO_Y=\OO_Q\oplus\OO_Q(-2,-2),\qquad
 \omega_Y\simeq\pi^*(\omega_Q\otimes L)\simeq\OO_Y.
\]
The first identity follows directly from the double-cover construction;
the second is its adjunction formula. Since $\pi$ is finite,
K\"unneth's formula gives
\[
 H^1(Y,\OO_Y)
 =H^1(Q,\OO_Q)\oplus H^1(Q,\OO_Q(-2,-2))=0.
\]
Indeed each summand in the latter K\"unneth expression contains
$H^0(\Pp^1,\OO(-2))=0$. Singularities of type $A_3$ are rational
singularities, and their minimal resolutions are crepant.
Thus resolution preserves this $H^1$ and the trivial canonical
sheaf. This is also the classical $(4,4)$ double-cover criterion
stated in \cite[Section 3, p.~319]{AngermullerBarth}.
The smooth resolution is therefore a K3 surface.
\end{proof}

For completeness, the local crepant calculation can be seen directly
on $ab=c^4$. Blowing up its origin leaves, in the $c$ chart, the
singularity $a_1b_1=c^2$; its next blow-up is smooth. The other
charts are already smooth. The canonical form $da\wedge dc/a$
becomes $da\wedge dc_1$ in the $a$ chart and
$da_1\wedge dc/a_1$ in the $c$ chart. The second blow-up has
the same calculation, with a nowhere vanishing form on each smooth
chart. The three exceptional curves form the usual $A_3$ chain.

Over $\mathbb C$, the topological Euler numbers give an independent
numerical check of the canonical calculation.
Since both branch components have genus one and meet in four points,
$e(B_1\cup B_2)=-4$. Hence $e(Y)=2e(Q)-e(B_1\cup B_2)=12$.
Resolving the four $A_3$ points adds $4\cdot3$, giving
$e(\widetilde Y)=24$.

\section{The common-parameter fiber product}

Fix a finite list $C_i,D_i\in k^\times$, $1\le i\le r$, with
$r\ge1$. For each, keep the map to the common $\tau$ line given
by \eqref{eq:affine}. Let $Z$ be the component dominating
$\Pp^1_\tau$ of the fiber product of these one-pole surfaces.

\begin{corollary}\label{cor:notrc}
The generic fiber of $Z\to\Pp^1_\tau$ is geometrically integral,
and $Z$ dominates each $Y_{C_i,D_i}$. No smooth projective model
of $Z$ is geometrically rationally connected.
\end{corollary}
\begin{proof}
Over an algebraic closure of $k(\tau)$, the $i$th generic fiber
is the double cover
\[
       q_i^2=D_i(s_i^2-A_i)(s_i^2+3A_i),\qquad
       A_i=C_i\frac{\tau^2-D_i}{\tau^2+3D_i}\ne0.
\]
Its four branch points are distinct, so its smooth projective
model is a geometrically integral genus-one curve. The product
of these curves is geometrically integral and projects
dominantly onto each factor. This proves the generic-fiber and
dominance assertions, and also the uniqueness of the component
dominating the parameter line.

After resolving the projection, a smooth projective model of $Z$
dominates the K3 surface $\widetilde Y_{C_i,D_i}$. Pullback of
its nonzero regular two-form is nonzero in characteristic zero.
Regular differential forms are birational invariants of smooth
proper varieties, so every smooth projective model of $Z$ has
a nonzero regular two-form. Such a variety cannot be rationally
connected, by the differential-form vanishing theorem
\cite[Proposition 2.5(b)]{Debarre}.
\end{proof}

\paragraph{Arithmetic scope.}
The coefficients in this result are fixed detector values and their
extracted square coefficients. The conic parameter remains variable.
Thus the obstruction is stronger than the statement that one
fixed-parameter fiber has genus one, but applies to this particular
fixed-input family. It does not assert an obstruction to rational
points, to a shared rational parameter, to rational subvarieties, or
to Zariski density. It does not by itself decide uniruledness of
the higher-dimensional fiber product. Varying auxiliary elliptic
indices changes the detector values and is a separate arithmetic
construction; the independent-index theorem of \cite{Guards}
remains applicable. Direct rational-connectedness results cannot
be applied to the smooth projective models in
Corollary~\ref{cor:notrc}.

\section{Exact algebraic verification}

The checker
\begin{center}\small
\texttt{negative\_results/computation/fixed\_input\_k3\_v42/check.py}
\end{center}
verifies the bihomogeneous branch equations, the smoothness determinants
in all four affine charts, the intersection identities, both local
normal forms and tangency orders, and the two blow-ups of $ab=c^4$.
It also checks the canonical-form transformations and the Euler and
cohomology dimension calculations. The global K3 and
rational-connectedness conclusions use the proofs and geometric
results stated above; they are not inferred from numerical examples.

\begin{thebibliography}{9}\raggedright
\bibitem{AngermullerBarth}
G.~Angerm\"uller and W.~Barth,
\emph{Elliptic fibres on Enriques surfaces},
Compositio Mathematica \textbf{47} (1982), no.~3, 317--332,
Section~3, p.~319,
\url{https://www.numdam.org/item/CM_1982__47_3_317_0.pdf}.
\bibitem{Debarre}
O.~Debarre, \emph{On rationality problems},
author's notes, 15 April 2024, Proposition~2.5(b),
\url{https://perso.imj-prg.fr/olivier-debarre/wp-content/uploads/sites/34/2024/04/On-rationality-problems.pdf}.
\bibitem{Guards}
\emph{Fixed pole fibers and independent detector indices},
project paper, version~42,
\texttt{papers/pole\_parameters/}\allowbreak
\texttt{common\_parameter\_guard\_reduction\_20261116\_v42.tex}.
\end{thebibliography}
\end{document}
